\documentclass[a4paper,12pt]{article}
\usepackage[utf8]{inputenc}
\usepackage[T1]{fontenc}
\usepackage[french]{babel}
\usepackage{geometry}
\geometry{left=1.8cm,right=1.8cm,top=1.3cm,bottom=1.6cm}

\usepackage{lmodern}
\usepackage{xcolor}
\definecolor{titleblue}{RGB}{0,51,140}
\definecolor{TITLEBLUE}{RGB}{0,51,140}

\usepackage{hyperref}
\hypersetup{
    colorlinks=true,
    urlcolor=titleblue,
    linkcolor=titleblue,
    citecolor=titleblue,
    pdfborder={0 0 0}
}

\usepackage{titlesec}
\usepackage{enumitem}

\pagestyle{empty}

\titleformat{\section}
  {\color{titleblue}\Large\bfseries\uppercase}{}{0em}{}
  [\color{titleblue}\titlerule]

\titlespacing\section{0pt}{8pt}{4pt}
\setlist[itemize]{leftmargin=*,topsep=2pt,itemsep=1pt}

\begin{document}

\begin{center}
    {\Huge\bfseries\color{titleblue} AISSAOUI Ismail}\\[4pt]
    {\Large Ingénieur IA \& Data Science}\\[6pt]

    {\normalsize
    Email : \href{mailto:ismail.aissaoui.pro@gmail.com}{ismail.aissaoui.pro@gmail.com} \quad
    Tél : +213 660 70 77 96 \\
    Portfolio : \href{https://ismail-aissaoui.vercel.app}{ismail-aissaoui.vercel.app} \quad
    Localisation : Chlef, Algérie \\
    LinkedIn : \href{https://linkedin.com/in/aissaoui-ismail-6a77a92a7}{linkedin.com/in/aissaoui-ismail-6a77a92a7} \quad
    GitHub : \href{https://github.com/i-aissaoui}{github.com/i-aissaoui}
    }
\end{center}

%===========================
\section{RÉSUMÉ PROFESSIONNEL}
Ingénieur IA \& Data Science spécialisé dans l'apprentissage de représentations latentes, les réseaux de graphes (GNN), et les modèles Transformers. Passionné par l'IA explicable (XAI) et la modélisation de relations de similarité complexes à partir de données multimodales. Candidat au doctorat pour l'aide au raisonnement clinique via la construction de graphes de patients et l'évaluation de la représentativité/atypicité.

%===========================
\section{COMPÉTENCES}
\setlength{\parskip}{-1pt}
{\normalsize
\textbf{Programmation}: Python, TypeScript, JavaScript, Java, SQL \\

\textbf{ML / DL}: PyTorch, TensorFlow, Keras, Scikit-learn, \textbf{Metric Learning}, \textbf{Autoencodeurs} \\

\textbf{Techniques IA}: GNN, Transformers, Embeddings Latents, \textbf{IA Explicable (XAI)}, Apprentissage Multimodal \\

\textbf{Optimisation}: PSO, Algorithme génétique, ACO, Recuit simulé, Optimisation bayésienne \\

\textbf{Web \& API}: React, Next.js, FastAPI, Node.js, TailwindCSS \\

\textbf{Data \& MLOps}: MLflow, Airflow, Kubeflow, DVC, Pandas, NumPy \\

\textbf{Bases de données}: PostgreSQL, MongoDB, MySQL, Neo4j, Cassandra \\

\textbf{Cloud \& DevOps}: Docker, Linux, Git, CI/CD, CUDA, AWS, Grafana \\

\textbf{Soft Skills}: Leadership technique, Communication, Gestion des parties prenantes
}

%===========================
\section{EXPÉRIENCE PROFESSIONNELLE}
\textbf{Stagiaire Ingénieur IA — Jumeau Numérique \& Maintenance Prédictive} \hfill 2026 \\
Sonatrach — Algérie \\
Développement d’un jumeau numérique pour turbines à gaz permettant une surveillance en temps réel et une maintenance prédictive. Implémentation d’une détection d’anomalies avancée et estimation de la RUL à l’aide de modèles de deep learning.

\textbf{Stagiaire en ingénierie réseau} \hfill Sep 2024 \\
Algérie Télécom RMC Center — Chlef, Algérie

%===========================
\section{FORMATION}
\textbf{École Nationale Supérieure d’Informatique (ESI-SBA)} \hfill 2021--2026 \\
Diplôme d’ingénieur et master en informatique \\
Spécialisation : Intelligence Artificielle \& Informatique

%===========================
\section{RÉALISATIONS CLÉS}
\begin{itemize}
    \item Développement de modèles Transformer (BERT, GPT, RAG) pour des applications multilingues réelles.
    \item Conception de plateformes d’IA explicable (XAI) pour la HR tech et les institutions publiques.
    \item Apprentissage de représentations sur graphes de 60\,000 nœuds (+24\% de précision en similarité).
    \item Pipelines ML complets : ingestion multimodale, validation, entraînement, et monitoring.
    \item Systèmes interactifs humains-dans-la-boucle réduisant les biais et le drift de 32\%.
\end{itemize}

%===========================
\section{PROJETS}
\setlength{\parskip}{6pt}
{\small

\textbf{Graph-Based Similarity Intelligence (Career Connect)} \hfill 2024--2025 \\
Modélisation de relations de similarité complexes entre profils via des \textbf{GNN}. Apprentissage de représentations latentes pour l'appariement sémantique performant et l'identification de profils atypiques. \\
\textit{Tech : PyTorch Geometric, SBERT, FastAPI — Concepts : GNN, Embeddings relationnels, Metric Learning}

\textbf{Recognizini — Latent Representation \& Feedback} \hfill 2025 \\
Génération et mise à jour d'embeddings latents avec boucle de feedback pour corriger le modèle et identifier des cas aberrants. Un parallèle direct avec le raisonnement clinique interactif. \\
\textit{Tech : OpenCV, PyTorch, FastAPI — Concepts : Embeddings, Détection d'anomalies, Label Propagation}

\textbf{Pipeline NLP multimodal multi-optimisation} \hfill 2025 \\
Architecture de classification optimisée via PSO et algorithm génétique. Gestion de la robustesse et stabilité des résultats de classification de textes complexes. \\
\textit{Tech : DistilBERT, PyTorch, FastAPI — Concepts : Modèles Multimodaux, Optimisation, Robustesse}

\textbf{Mini-GPT (Decoder-only)} \hfill 2025 \\
Implémentation expérimentale d’un GPT compact pour comprendre la génération séquentielle et la stabilité des LLM. Prise en main avancée des architectures d'attention. \\
\textit{Tech : PyTorch, CUDA, AdamW — Concepts : Transformers, Modèles de langage, Représentation}

\textbf{North African Sentiment BERT} \hfill 2025 \\
Adaptation LoRA de BERT pour l’analyse de données textuelles fortement bruitées et non standardisées. \\
\textit{Tech : PyTorch, LoRA, MLflow — Concepts : Fine-tuning, LLM, XAI}

\textbf{Geo-RAG \& Augmented Generation} \hfill 2025 \\
Système RAG intégrant de la donnée structurée (graphes, geo) et non structurée (texte) pour des réponses explicables. \\
\textit{Tech : ChromaDB, LangChain — Concepts : RAG, Text Mining, Explicabilité}

\textbf{Elliptic++ : Détection d'anomalies sur Graphes} \hfill 2025 \\
Identification de "cas" frauduleux dans des réseaux massifs de transactions (200k+ nœuds) via auto-encodeurs sur graphes. \\
\textit{Tech : PyTorch Geometric, NetworkX — Concepts : GNN, Détection d'atypicités}
}



%===========================
\section{LANGUES}
Arabe : Langue maternelle \quad|\quad Anglais : Avancé \quad|\quad Français : Conversationnel

\end{document}
